%%%%%%%%%%%%%%%%%%%%%%%%%%%%%%%%%%%%%%%%%%%%%%%%%%%%%%%%%%%%
% 7.2 RANDOM VARIABLES
% The Composition of Advanced Mathematics
%%%%%%%%%%%%%%%%%%%%%%%%%%%%%%%%%%%%%%%%%%%%%%%%%%%%%%%%%%%%

\section{Random Variables}
\label{sec:random-variables}

\prerequisite{
Elementary probability, sample spaces, events, functions, intervals,
sets, inequalities, and basic algebra.
}

\learningobjectives{
By the end of this section, the reader should be able to:
\begin{itemize}
    \item define a random variable as a function on a sample space;
    \item distinguish outcomes from numerical random-variable values;
    \item identify the support of a random variable;
    \item distinguish discrete, continuous, and mixed random variables;
    \item determine probabilities of events involving random variables;
    \item define and use probability mass functions;
    \item understand probability density functions at an introductory
          level;
    \item define and interpret cumulative distribution functions;
    \item use the fundamental properties of cumulative distribution
          functions;
    \item recover interval probabilities from a distribution function;
    \item understand point masses and jumps of a distribution function;
    \item recognize degenerate random variables;
    \item define indicator random variables;
    \item use indicators to encode events numerically;
    \item construct random variables from common experiments;
    \item understand transformations of random variables;
    \item determine transformed supports;
    \item distinguish one-to-one and many-to-one transformations;
    \item understand location and scale transformations;
    \item define quantiles and medians using the distribution function;
    \item recognize random vectors as extensions to multiple random
          variables;
    \item understand how random variables connect probability spaces
          with numerical analysis.
\end{itemize}
}

%%%%%%%%%%%%%%%%%%%%%%%%%%%%%%%%%%%%%%%%%%%%%%%%%%%%%%%%%%%%
\subsection{From Outcomes to Numbers}
\label{subsec:from-outcomes-to-numbers}
%%%%%%%%%%%%%%%%%%%%%%%%%%%%%%%%%%%%%%%%%%%%%%%%%%%%%%%%%%%%

A sample space may contain outcomes that are not themselves numerical.

For example, two coin tosses produce

\[
\Omega
=
\{HH,HT,TH,TT\}.
\]

Suppose we care only about the number of heads.

We can assign

\[
HH\longmapsto2,
\]

\[
HT\longmapsto1,
\]

\[
TH\longmapsto1,
\]

and

\[
TT\longmapsto0.
\]

This assignment converts the outcomes of the experiment into numbers.

That function is a random variable.

%%%%%%%%%%%%%%%%%%%%%%%%%%%%%%%%%%%%%%%%%%%%%%%%%%%%%%%%%%%%
\subsection{Definition of a Random Variable}
\label{subsec:random-variable-definition}
%%%%%%%%%%%%%%%%%%%%%%%%%%%%%%%%%%%%%%%%%%%%%%%%%%%%%%%%%%%%

\begin{definitionbox}{Random Variable}
A random variable is a function

\[
\boxed{
X:\Omega\to\mathbb{R}
}
\]

that assigns a real number \(X(\omega)\) to each outcome

\[
\omega\in\Omega.
\]
\end{definitionbox}

The notation

\[
X(\omega)
\]

is read

\[
\text{``\(X\) evaluated at omega.''}
\]

The randomness comes from the uncertain outcome \(\omega\).

The function \(X\) itself is mathematically fixed.

%%%%%%%%%%%%%%%%%%%%%%%%%%%%%%%%%%%%%%%%%%%%%%%%%%%%%%%%%%%%
\subsection{Random Variables Are Functions}
\label{subsec:random-variables-functions}
%%%%%%%%%%%%%%%%%%%%%%%%%%%%%%%%%%%%%%%%%%%%%%%%%%%%%%%%%%%%

It is important to remember:

\[
\boxed{
\text{a random variable is a function, not an event}.
}
\]

The probability model determines which outcome

\[
\omega
\]

occurs.

The random variable then assigns that outcome a numerical value.

Thus,

\[
\boxed{
\omega
\longrightarrow
X(\omega).
}
\]

This permits numerical tools such as algebra, calculus, and analysis to
be applied to probability.

%%%%%%%%%%%%%%%%%%%%%%%%%%%%%%%%%%%%%%%%%%%%%%%%%%%%%%%%%%%%
\subsection{Worked Example: Number of Heads}
\label{subsec:worked-rv-heads}
%%%%%%%%%%%%%%%%%%%%%%%%%%%%%%%%%%%%%%%%%%%%%%%%%%%%%%%%%%%%

\begin{workedexamplebox}{Two Coin Tosses}

Let

\[
\Omega
=
\{HH,HT,TH,TT\}.
\]

Define

\[
X
=
\text{number of heads}.
\]

Then

\[
X(HH)=2,
\]

\[
X(HT)=1,
\]

\[
X(TH)=1,
\]

and

\[
X(TT)=0.
\]

Therefore the possible values of \(X\) are

\[
\{0,1,2\}.
\]

\begin{answerbox}
\[
\boxed{
X(\Omega)=\{0,1,2\}.
}
\]
\end{answerbox}

\end{workedexamplebox}

%%%%%%%%%%%%%%%%%%%%%%%%%%%%%%%%%%%%%%%%%%%%%%%%%%%%%%%%%%%%
\subsection{Outcomes Versus Random-Variable Values}
\label{subsec:outcomes-versus-rv-values}
%%%%%%%%%%%%%%%%%%%%%%%%%%%%%%%%%%%%%%%%%%%%%%%%%%%%%%%%%%%%

Several different outcomes may produce the same random-variable value.

In the previous example,

\[
HT
\]

and

\[
TH
\]

are distinct outcomes, but

\[
X(HT)=X(TH)=1.
\]

Thus,

\[
\boxed{
\text{outcome}
\neq
\text{random-variable value}.
}
\]

The random variable may discard information about the original
experiment.

%%%%%%%%%%%%%%%%%%%%%%%%%%%%%%%%%%%%%%%%%%%%%%%%%%%%%%%%%%%%
\subsection{Events Defined by Random Variables}
\label{subsec:events-defined-rv}
%%%%%%%%%%%%%%%%%%%%%%%%%%%%%%%%%%%%%%%%%%%%%%%%%%%%%%%%%%%%

Statements about random variables correspond to events in the original
sample space.

For example,

\[
\{X=1\}
\]

means

\[
\boxed{
\{\omega\in\Omega:X(\omega)=1\}.
}
\]

Likewise,

\[
\{X\leq x\}
\]

means

\[
\boxed{
\{\omega\in\Omega:X(\omega)\leq x\}.
}
\]

Thus probabilities such as

\[
P(X=1)
\]

are shorthand for probabilities of events.

%%%%%%%%%%%%%%%%%%%%%%%%%%%%%%%%%%%%%%%%%%%%%%%%%%%%%%%%%%%%
\subsection{Preimages of Sets}
\label{subsec:rv-preimages}
%%%%%%%%%%%%%%%%%%%%%%%%%%%%%%%%%%%%%%%%%%%%%%%%%%%%%%%%%%%%

If

\[
A\subseteq\mathbb{R},
\]

then

\[
\{X\in A\}
\]

means the preimage

\[
\boxed{
X^{-1}(A)
=
\{\omega\in\Omega:X(\omega)\in A\}.
}
\]

Therefore,

\[
\boxed{
P(X\in A)
=
P\bigl(X^{-1}(A)\bigr).
}
\]

Random variables transfer probability from the sample space to subsets
of the real line.

%%%%%%%%%%%%%%%%%%%%%%%%%%%%%%%%%%%%%%%%%%%%%%%%%%%%%%%%%%%%
\subsection{Distribution of a Random Variable}
\label{subsec:distribution-random-variable}
%%%%%%%%%%%%%%%%%%%%%%%%%%%%%%%%%%%%%%%%%%%%%%%%%%%%%%%%%%%%

The distribution of \(X\) describes how probability is assigned to its
possible numerical values.

Informally,

\[
\boxed{
\text{distribution of }X
=
\text{probability behavior of the values of }X.
}
\]

Two different experiments can produce random variables with the same
distribution.

In that case, the underlying sample spaces may differ even though the
numerical probability laws are identical.

%%%%%%%%%%%%%%%%%%%%%%%%%%%%%%%%%%%%%%%%%%%%%%%%%%%%%%%%%%%%
\subsection{Support}
\label{subsec:random-variable-support}
%%%%%%%%%%%%%%%%%%%%%%%%%%%%%%%%%%%%%%%%%%%%%%%%%%%%%%%%%%%%

\begin{definitionbox}{Support}
The support of a random variable is, informally, the set of values where
its probability distribution is concentrated.
\end{definitionbox}

For a discrete random variable, the support may be written

\[
\boxed{
\operatorname{supp}(X)
=
\{x:P(X=x)>0\}.
}
\]

For continuous distributions, the precise definition uses topological
closure and will be refined later.

%%%%%%%%%%%%%%%%%%%%%%%%%%%%%%%%%%%%%%%%%%%%%%%%%%%%%%%%%%%%
\subsection{Worked Example: Die Random Variable}
\label{subsec:worked-die-rv}
%%%%%%%%%%%%%%%%%%%%%%%%%%%%%%%%%%%%%%%%%%%%%%%%%%%%%%%%%%%%

\begin{workedexamplebox}{Square of a Die Roll}

Roll a fair six-sided die and let

\[
X
=
(\text{number rolled})^2.
\]

Then the possible values are

\[
1,4,9,16,25,36.
\]

Since each original die result has probability

\[
\frac16,
\]

we have

\[
P(X=1)
=
P(X=4)
=
\cdots
=
P(X=36)
=
\frac16.
\]

\begin{answerbox}
\[
\boxed{
\operatorname{supp}(X)
=
\{1,4,9,16,25,36\}.
}
\]
\end{answerbox}

\end{workedexamplebox}

%%%%%%%%%%%%%%%%%%%%%%%%%%%%%%%%%%%%%%%%%%%%%%%%%%%%%%%%%%%%
\subsection{Discrete Random Variables}
\label{subsec:discrete-random-variables}
%%%%%%%%%%%%%%%%%%%%%%%%%%%%%%%%%%%%%%%%%%%%%%%%%%%%%%%%%%%%

\begin{definitionbox}{Discrete Random Variable}
A random variable is discrete if its possible values form a finite or
countably infinite set.
\end{definitionbox}

Examples include:

\[
\text{number of heads in \(10\) tosses},
\]

\[
\text{number of customers arriving},
\]

\[
\text{number of defective components},
\]

and

\[
\text{number of trials before the first success}.
\]

%%%%%%%%%%%%%%%%%%%%%%%%%%%%%%%%%%%%%%%%%%%%%%%%%%%%%%%%%%%%
\subsection{Probability Mass Function}
\label{subsec:probability-mass-function}
%%%%%%%%%%%%%%%%%%%%%%%%%%%%%%%%%%%%%%%%%%%%%%%%%%%%%%%%%%%%

\begin{definitionbox}{Probability Mass Function}
For a discrete random variable \(X\), the probability mass function,
abbreviated PMF, is

\[
\boxed{
p_X(x)
=
P(X=x).
}
\]
\end{definitionbox}

The subscript \(X\) indicates that the function belongs to random
variable \(X\).

The PMF completely determines the distribution of a discrete random
variable.

%%%%%%%%%%%%%%%%%%%%%%%%%%%%%%%%%%%%%%%%%%%%%%%%%%%%%%%%%%%%
\subsection{Properties of a PMF}
\label{subsec:pmf-properties}
%%%%%%%%%%%%%%%%%%%%%%%%%%%%%%%%%%%%%%%%%%%%%%%%%%%%%%%%%%%%

A probability mass function satisfies

\[
\boxed{
p_X(x)\geq0
}
\]

for every \(x\), and

\[
\boxed{
\sum_x p_X(x)=1.
}
\]

The sum is taken over all possible values of \(X\).

For any set \(A\),

\[
\boxed{
P(X\in A)
=
\sum_{x\in A}p_X(x).
}
\]

%%%%%%%%%%%%%%%%%%%%%%%%%%%%%%%%%%%%%%%%%%%%%%%%%%%%%%%%%%%%
\subsection{Worked Example: PMF of the Number of Heads}
\label{subsec:worked-pmf-heads}
%%%%%%%%%%%%%%%%%%%%%%%%%%%%%%%%%%%%%%%%%%%%%%%%%%%%%%%%%%%%

\begin{workedexamplebox}{Two Fair Coin Tosses}

Let

\[
X
=
\text{number of heads in two fair coin tosses}.
\]

The four equally likely outcomes are

\[
HH,HT,TH,TT.
\]

Therefore,

\[
P(X=0)=\frac14,
\]

\[
P(X=1)=\frac24=\frac12,
\]

and

\[
P(X=2)=\frac14.
\]

Thus,

\begin{answerbox}
\[
\boxed{
p_X(x)
=
\begin{cases}
\dfrac14,&x=0,\\[2mm]
\dfrac12,&x=1,\\[2mm]
\dfrac14,&x=2,\\[2mm]
0,&\text{otherwise}.
\end{cases}
}
\]
\end{answerbox}

\end{workedexamplebox}

%%%%%%%%%%%%%%%%%%%%%%%%%%%%%%%%%%%%%%%%%%%%%%%%%%%%%%%%%%%%
\subsection{PMFs Need Not Be Uniform}
\label{subsec:nonuniform-pmf}
%%%%%%%%%%%%%%%%%%%%%%%%%%%%%%%%%%%%%%%%%%%%%%%%%%%%%%%%%%%%

A discrete random variable may have equally spaced numerical values but
unequal probabilities.

For example, if

\[
p_X(0)=0.1,
\qquad
p_X(1)=0.3,
\qquad
p_X(2)=0.6,
\]

then

\[
\boxed{
0.1+0.3+0.6=1,
}
\]

so this defines a valid PMF.

The numerical spacing of the values does not determine their
probabilities.

%%%%%%%%%%%%%%%%%%%%%%%%%%%%%%%%%%%%%%%%%%%%%%%%%%%%%%%%%%%%
\subsection{Countably Infinite Discrete Variables}
\label{subsec:countably-infinite-rv}
%%%%%%%%%%%%%%%%%%%%%%%%%%%%%%%%%%%%%%%%%%%%%%%%%%%%%%%%%%%%

Discrete random variables need not have finite support.

For example, repeatedly toss a coin until the first head appears.

Let

\[
X
=
\text{number of tosses required}.
\]

Then

\[
\boxed{
X\in\{1,2,3,\ldots\}.
}
\]

The support is countably infinite.

Such random variables will be developed further in the discrete
distribution section.

%%%%%%%%%%%%%%%%%%%%%%%%%%%%%%%%%%%%%%%%%%%%%%%%%%%%%%%%%%%%
\subsection{Continuous Random Variables}
\label{subsec:continuous-random-variables}
%%%%%%%%%%%%%%%%%%%%%%%%%%%%%%%%%%%%%%%%%%%%%%%%%%%%%%%%%%%%

\begin{definitionbox}{Continuous Random Variable}
A continuous random variable is one whose distribution can be described
by a probability density function over intervals of real numbers.
\end{definitionbox}

Typical examples include:

\[
\text{time},
\qquad
\text{length},
\qquad
\text{temperature},
\qquad
\text{mass},
\qquad
\text{voltage}.
\]

Continuous random variables generally take values in intervals rather
than countable sets.

%%%%%%%%%%%%%%%%%%%%%%%%%%%%%%%%%%%%%%%%%%%%%%%%%%%%%%%%%%%%
\subsection{Probability Density Function Preview}
\label{subsec:pdf-preview}
%%%%%%%%%%%%%%%%%%%%%%%%%%%%%%%%%%%%%%%%%%%%%%%%%%%%%%%%%%%%

For a continuous random variable \(X\), a probability density function,
abbreviated PDF, is a function

\[
f_X(x)
\]

such that

\[
\boxed{
P(a<X\leq b)
=
\int_a^b f_X(x)\,dx.
}
\]

A valid density satisfies

\[
\boxed{
f_X(x)\geq0
}
\]

and

\[
\boxed{
\int_{-\infty}^{\infty}f_X(x)\,dx=1.
}
\]

Continuous distributions will be studied systematically in
Section~\ref{sec:continuous-distributions}.

%%%%%%%%%%%%%%%%%%%%%%%%%%%%%%%%%%%%%%%%%%%%%%%%%%%%%%%%%%%%
\subsection{Density Is Not Probability}
\label{subsec:density-not-probability}
%%%%%%%%%%%%%%%%%%%%%%%%%%%%%%%%%%%%%%%%%%%%%%%%%%%%%%%%%%%%

For a continuous random variable,

\[
f_X(x)
\]

is a density value.

It is not generally equal to

\[
P(X=x).
\]

In fact, for a continuous random variable,

\[
\boxed{
P(X=x)=0
}
\]

for every individual real number \(x\).

Probabilities are obtained from areas under the density curve.

%%%%%%%%%%%%%%%%%%%%%%%%%%%%%%%%%%%%%%%%%%%%%%%%%%%%%%%%%%%%
\subsection{Continuous Interval Probabilities}
\label{subsec:continuous-interval-probabilities}
%%%%%%%%%%%%%%%%%%%%%%%%%%%%%%%%%%%%%%%%%%%%%%%%%%%%%%%%%%%%

Because single points have probability zero,

\[
P(a<X<b),
\]

\[
P(a\leq X<b),
\]

\[
P(a<X\leq b),
\]

and

\[
P(a\leq X\leq b)
\]

are equal for a continuous random variable.

Thus endpoint inclusion does not affect continuous interval
probabilities.

%%%%%%%%%%%%%%%%%%%%%%%%%%%%%%%%%%%%%%%%%%%%%%%%%%%%%%%%%%%%
\subsection{Worked Example: Uniform Interval}
\label{subsec:worked-uniform-rv-preview}
%%%%%%%%%%%%%%%%%%%%%%%%%%%%%%%%%%%%%%%%%%%%%%%%%%%%%%%%%%%%

\begin{workedexamplebox}{Uniform Random Point}

Suppose \(X\) is chosen uniformly from

\[
[0,10].
\]

Its density is

\[
f_X(x)
=
\frac1{10}
\]

for

\[
0\leq x\leq10.
\]

Find

\[
P(2\leq X\leq5).
\]

Then

\[
P(2\leq X\leq5)
=
\int_2^5\frac1{10}\,dx.
\]

Therefore,

\begin{answerbox}
\[
\boxed{
P(2\leq X\leq5)
=
\frac3{10}.
}
\]
\end{answerbox}

\end{workedexamplebox}

%%%%%%%%%%%%%%%%%%%%%%%%%%%%%%%%%%%%%%%%%%%%%%%%%%%%%%%%%%%%
\subsection{Mixed Random Variables}
\label{subsec:mixed-random-variables}
%%%%%%%%%%%%%%%%%%%%%%%%%%%%%%%%%%%%%%%%%%%%%%%%%%%%%%%%%%%%

A distribution may contain both:

\[
\text{discrete point masses}
\]

and

\[
\text{continuous probability}.
\]

Such a variable is called mixed.

For example, a machine lifetime might satisfy

\[
P(X=0)>0
\]

if some machines fail immediately, while surviving machines have a
continuous lifetime distribution.

Thus not every random variable is purely discrete or purely continuous.

%%%%%%%%%%%%%%%%%%%%%%%%%%%%%%%%%%%%%%%%%%%%%%%%%%%%%%%%%%%%
\subsection{Degenerate Random Variables}
\label{subsec:degenerate-random-variables}
%%%%%%%%%%%%%%%%%%%%%%%%%%%%%%%%%%%%%%%%%%%%%%%%%%%%%%%%%%%%

A random variable may be constant.

Suppose

\[
X(\omega)=c
\]

for every outcome.

Then

\[
\boxed{
P(X=c)=1.
}
\]

Such a random variable is called degenerate.

Its entire distribution is concentrated at one point.

%%%%%%%%%%%%%%%%%%%%%%%%%%%%%%%%%%%%%%%%%%%%%%%%%%%%%%%%%%%%
\subsection{Cumulative Distribution Function}
\label{subsec:cumulative-distribution-function}
%%%%%%%%%%%%%%%%%%%%%%%%%%%%%%%%%%%%%%%%%%%%%%%%%%%%%%%%%%%%

Every real-valued random variable has a cumulative distribution
function.

\begin{definitionbox}{Cumulative Distribution Function}
The cumulative distribution function, abbreviated CDF, of \(X\) is

\[
\boxed{
F_X(x)
=
P(X\leq x).
}
\]
\end{definitionbox}

The notation is read

\[
\text{``\(F_X\) of \(x\).''}
\]

The CDF applies to discrete, continuous, and mixed random variables.

%%%%%%%%%%%%%%%%%%%%%%%%%%%%%%%%%%%%%%%%%%%%%%%%%%%%%%%%%%%%
\subsection{Meaning of the CDF}
\label{subsec:cdf-meaning}
%%%%%%%%%%%%%%%%%%%%%%%%%%%%%%%%%%%%%%%%%%%%%%%%%%%%%%%%%%%%

For a given value \(x\),

\[
F_X(x)
\]

is the total probability accumulated at or below \(x\).

Thus the CDF records probability from

\[
-\infty
\]

up to the chosen value.

Informally,

\[
\boxed{
F_X(x)
=
\text{probability accumulated through }x.
}
\]

%%%%%%%%%%%%%%%%%%%%%%%%%%%%%%%%%%%%%%%%%%%%%%%%%%%%%%%%%%%%
\subsection{CDF Properties}
\label{subsec:cdf-properties}
%%%%%%%%%%%%%%%%%%%%%%%%%%%%%%%%%%%%%%%%%%%%%%%%%%%%%%%%%%%%

Every CDF satisfies:

\[
\boxed{
0\leq F_X(x)\leq1;
}
\]

\[
\boxed{
F_X(x)
\text{ is nondecreasing};
}
\]

\[
\boxed{
\lim_{x\to-\infty}F_X(x)=0;
}
\]

\[
\boxed{
\lim_{x\to\infty}F_X(x)=1;
}
\]

and \(F_X\) is right-continuous.

%%%%%%%%%%%%%%%%%%%%%%%%%%%%%%%%%%%%%%%%%%%%%%%%%%%%%%%%%%%%
\subsection{Why a CDF Is Nondecreasing}
\label{subsec:cdf-nondecreasing}
%%%%%%%%%%%%%%%%%%%%%%%%%%%%%%%%%%%%%%%%%%%%%%%%%%%%%%%%%%%%

Suppose

\[
x_1<x_2.
\]

Then

\[
\{X\leq x_1\}
\subseteq
\{X\leq x_2\}.
\]

By monotonicity of probability,

\[
P(X\leq x_1)
\leq
P(X\leq x_2).
\]

Therefore,

\[
\boxed{
F_X(x_1)\leq F_X(x_2).
}
\]

%%%%%%%%%%%%%%%%%%%%%%%%%%%%%%%%%%%%%%%%%%%%%%%%%%%%%%%%%%%%
\subsection{CDF of a Discrete Variable}
\label{subsec:cdf-discrete-variable}
%%%%%%%%%%%%%%%%%%%%%%%%%%%%%%%%%%%%%%%%%%%%%%%%%%%%%%%%%%%%

For a discrete variable,

\[
\boxed{
F_X(x)
=
\sum_{t\leq x}p_X(t).
}
\]

Thus the CDF is a step function.

Each positive probability mass creates a jump.

%%%%%%%%%%%%%%%%%%%%%%%%%%%%%%%%%%%%%%%%%%%%%%%%%%%%%%%%%%%%
\subsection{Worked Example: Discrete CDF}
\label{subsec:worked-discrete-cdf}
%%%%%%%%%%%%%%%%%%%%%%%%%%%%%%%%%%%%%%%%%%%%%%%%%%%%%%%%%%%%

\begin{workedexamplebox}{CDF for Two Coin Tosses}

Let \(X\) be the number of heads in two fair tosses.

Recall

\[
P(X=0)=\frac14,
\]

\[
P(X=1)=\frac12,
\]

and

\[
P(X=2)=\frac14.
\]

Therefore,

\[
F_X(x)
=
\begin{cases}
0,&x<0,\\[1mm]
\dfrac14,&0\leq x<1,\\[2mm]
\dfrac34,&1\leq x<2,\\[2mm]
1,&x\geq2.
\end{cases}
\]

\begin{answerbox}
\[
\boxed{
F_X(x)
=
\begin{cases}
0,&x<0,\\
\dfrac14,&0\leq x<1,\\
\dfrac34,&1\leq x<2,\\
1,&x\geq2.
\end{cases}
}
\]
\end{answerbox}

\end{workedexamplebox}

%%%%%%%%%%%%%%%%%%%%%%%%%%%%%%%%%%%%%%%%%%%%%%%%%%%%%%%%%%%%
\subsection{Jumps of the CDF}
\label{subsec:cdf-jumps}
%%%%%%%%%%%%%%%%%%%%%%%%%%%%%%%%%%%%%%%%%%%%%%%%%%%%%%%%%%%%

For any random variable,

\[
\boxed{
P(X=x)
=
F_X(x)-F_X(x^-),
}
\]

where

\[
F_X(x^-)
=
\lim_{t\uparrow x}F_X(t).
\]

The notation

\[
t\uparrow x
\]

means that \(t\) approaches \(x\) from below.

Thus the size of the jump at \(x\) equals the probability mass at
\(x\).

%%%%%%%%%%%%%%%%%%%%%%%%%%%%%%%%%%%%%%%%%%%%%%%%%%%%%%%%%%%%
\subsection{Continuous CDFs}
\label{subsec:continuous-cdfs}
%%%%%%%%%%%%%%%%%%%%%%%%%%%%%%%%%%%%%%%%%%%%%%%%%%%%%%%%%%%%

If \(X\) has density \(f_X\), then

\[
\boxed{
F_X(x)
=
\int_{-\infty}^{x}
f_X(t)\,dt.
}
\]

Under suitable regularity,

\[
\boxed{
F_X'(x)=f_X(x).
}
\]

Thus:

\[
\boxed{
\text{CDF}
\longleftrightarrow
\text{accumulated density}.
}
\]

%%%%%%%%%%%%%%%%%%%%%%%%%%%%%%%%%%%%%%%%%%%%%%%%%%%%%%%%%%%%
\subsection{Recovering Probabilities from the CDF}
\label{subsec:recover-probabilities-cdf}
%%%%%%%%%%%%%%%%%%%%%%%%%%%%%%%%%%%%%%%%%%%%%%%%%%%%%%%%%%%%

The CDF completely determines the distribution.

For example,

\[
\boxed{
P(X>x)
=
1-F_X(x).
}
\]

Also,

\[
\boxed{
P(a<X\leq b)
=
F_X(b)-F_X(a).
}
\]

More generally,

\[
P(a\leq X\leq b)
=
F_X(b)-F_X(a^-).
\]

%%%%%%%%%%%%%%%%%%%%%%%%%%%%%%%%%%%%%%%%%%%%%%%%%%%%%%%%%%%%
\subsection{Worked Example: Interval Probability from a CDF}
\label{subsec:worked-cdf-interval}
%%%%%%%%%%%%%%%%%%%%%%%%%%%%%%%%%%%%%%%%%%%%%%%%%%%%%%%%%%%%

\begin{workedexamplebox}{Using a Distribution Function}

Suppose

\[
F_X(2)=0.35
\]

and

\[
F_X(5)=0.80.
\]

Then

\[
P(2<X\leq5)
=
F_X(5)-F_X(2).
\]

Therefore,

\begin{answerbox}
\[
\boxed{
P(2<X\leq5)=0.45.
}
\]
\end{answerbox}

\end{workedexamplebox}

%%%%%%%%%%%%%%%%%%%%%%%%%%%%%%%%%%%%%%%%%%%%%%%%%%%%%%%%%%%%
\subsection{CDF and Complement Probabilities}
\label{subsec:cdf-complements}
%%%%%%%%%%%%%%%%%%%%%%%%%%%%%%%%%%%%%%%%%%%%%%%%%%%%%%%%%%%%

Since

\[
F_X(x)=P(X\leq x),
\]

we obtain

\[
\boxed{
P(X>x)=1-F_X(x).
}
\]

For strict lower inequalities,

\[
P(X\geq x)
=
1-P(X<x).
\]

If \(X\) is continuous,

\[
P(X<x)=P(X\leq x),
\]

so

\[
\boxed{
P(X\geq x)=1-F_X(x).
}
\]

For discrete variables, endpoint masses must be handled carefully.

%%%%%%%%%%%%%%%%%%%%%%%%%%%%%%%%%%%%%%%%%%%%%%%%%%%%%%%%%%%%
\subsection{PMF from the CDF}
\label{subsec:pmf-from-cdf}
%%%%%%%%%%%%%%%%%%%%%%%%%%%%%%%%%%%%%%%%%%%%%%%%%%%%%%%%%%%%

For a discrete random variable,

\[
p_X(x)
=
P(X=x).
\]

This is exactly the jump of the CDF:

\[
\boxed{
p_X(x)
=
F_X(x)-F_X(x^-).
}
\]

Thus:

\[
\boxed{
\text{PMF}
=
\text{jump sizes of the CDF}.
}
\]

%%%%%%%%%%%%%%%%%%%%%%%%%%%%%%%%%%%%%%%%%%%%%%%%%%%%%%%%%%%%
\subsection{PDF from the CDF}
\label{subsec:pdf-from-cdf}
%%%%%%%%%%%%%%%%%%%%%%%%%%%%%%%%%%%%%%%%%%%%%%%%%%%%%%%%%%%%

If the CDF is differentiable and the distribution is absolutely
continuous, then

\[
\boxed{
f_X(x)=F_X'(x).
}
\]

Conversely,

\[
\boxed{
F_X(x)
=
\int_{-\infty}^{x}f_X(t)\,dt.
}
\]

The exact measure-theoretic conditions will be developed later.

%%%%%%%%%%%%%%%%%%%%%%%%%%%%%%%%%%%%%%%%%%%%%%%%%%%%%%%%%%%%
\subsection{Indicator Random Variables}
\label{subsec:indicator-random-variables-probability}
%%%%%%%%%%%%%%%%%%%%%%%%%%%%%%%%%%%%%%%%%%%%%%%%%%%%%%%%%%%%

An event can be converted into a random variable.

\begin{definitionbox}{Indicator Random Variable}
For an event \(A\), its indicator is

\[
\boxed{
\mathbf{1}_A(\omega)
=
\begin{cases}
1,&\omega\in A,\\
0,&\omega\notin A.
\end{cases}
}
\]
\end{definitionbox}

Thus,

\[
\boxed{
\mathbf{1}_A
=
1
\iff
A\text{ occurs}.
}
\]

%%%%%%%%%%%%%%%%%%%%%%%%%%%%%%%%%%%%%%%%%%%%%%%%%%%%%%%%%%%%
\subsection{Distribution of an Indicator}
\label{subsec:indicator-distribution}
%%%%%%%%%%%%%%%%%%%%%%%%%%%%%%%%%%%%%%%%%%%%%%%%%%%%%%%%%%%%

If

\[
P(A)=p,
\]

then

\[
P(\mathbf{1}_A=1)=p
\]

and

\[
P(\mathbf{1}_A=0)=1-p.
\]

Thus an indicator has the two-point distribution

\[
\boxed{
p_{\mathbf{1}_A}(x)
=
\begin{cases}
1-p,&x=0,\\
p,&x=1,\\
0,&\text{otherwise}.
\end{cases}
}
\]

This is the Bernoulli distribution.

%%%%%%%%%%%%%%%%%%%%%%%%%%%%%%%%%%%%%%%%%%%%%%%%%%%%%%%%%%%%
\subsection{Indicators as Counting Tools}
\label{subsec:indicators-counting-tools}
%%%%%%%%%%%%%%%%%%%%%%%%%%%%%%%%%%%%%%%%%%%%%%%%%%%%%%%%%%%%

Suppose events

\[
A_1,\ldots,A_n
\]

represent occurrences we want to count.

Define

\[
X
=
\mathbf{1}_{A_1}
+\cdots+
\mathbf{1}_{A_n}.
\]

Then \(X\) counts how many of the events occur.

Thus,

\[
\boxed{
\text{count}
=
\text{sum of indicators}.
}
\]

This idea becomes fundamental when studying expectation.

%%%%%%%%%%%%%%%%%%%%%%%%%%%%%%%%%%%%%%%%%%%%%%%%%%%%%%%%%%%%
\subsection{Worked Example: Counting Sixes}
\label{subsec:worked-indicator-sixes}
%%%%%%%%%%%%%%%%%%%%%%%%%%%%%%%%%%%%%%%%%%%%%%%%%%%%%%%%%%%%

\begin{workedexamplebox}{Indicators for Repeated Die Rolls}

Roll a die \(5\) times.

For roll \(i\), define

\[
I_i
=
\begin{cases}
1,&\text{roll \(i\) is a six},\\
0,&\text{otherwise}.
\end{cases}
\]

Then

\[
X
=
I_1+I_2+I_3+I_4+I_5
\]

is exactly the number of sixes rolled.

\begin{answerbox}
\[
\boxed{
X=\sum_{i=1}^{5}I_i.
}
\]
\end{answerbox}

\end{workedexamplebox}

%%%%%%%%%%%%%%%%%%%%%%%%%%%%%%%%%%%%%%%%%%%%%%%%%%%%%%%%%%%%
\subsection{Functions of Random Variables}
\label{subsec:functions-of-random-variables}
%%%%%%%%%%%%%%%%%%%%%%%%%%%%%%%%%%%%%%%%%%%%%%%%%%%%%%%%%%%%

If \(X\) is a random variable and

\[
g:\mathbb{R}\to\mathbb{R},
\]

then

\[
\boxed{
Y=g(X)
}
\]

defines another random variable.

More explicitly,

\[
Y(\omega)
=
g(X(\omega)).
\]

Thus random variables can be transformed using ordinary functions.

%%%%%%%%%%%%%%%%%%%%%%%%%%%%%%%%%%%%%%%%%%%%%%%%%%%%%%%%%%%%
\subsection{Worked Example: Squaring a Random Variable}
\label{subsec:worked-square-random-variable}
%%%%%%%%%%%%%%%%%%%%%%%%%%%%%%%%%%%%%%%%%%%%%%%%%%%%%%%%%%%%

\begin{workedexamplebox}{Transforming a Die Roll}

Let \(X\) be the number appearing on a fair die.

Define

\[
Y=X^2.
\]

Then

\[
Y
\in
\{1,4,9,16,25,36\}.
\]

Since \(x\mapsto x^2\) is one-to-one on

\[
\{1,2,3,4,5,6\},
\]

each value still has probability

\[
\frac16.
\]

\begin{answerbox}
\[
\boxed{
P(Y=y)=\frac16
}
\]

for

\[
y\in\{1,4,9,16,25,36\}.
\]
\end{answerbox}

\end{workedexamplebox}

%%%%%%%%%%%%%%%%%%%%%%%%%%%%%%%%%%%%%%%%%%%%%%%%%%%%%%%%%%%%
\subsection{Many-to-One Transformations}
\label{subsec:many-to-one-transformations}
%%%%%%%%%%%%%%%%%%%%%%%%%%%%%%%%%%%%%%%%%%%%%%%%%%%%%%%%%%%%

A transformation may map several values of \(X\) to the same value of
\(Y\).

Suppose

\[
Y=X^2
\]

and \(X\) can take both positive and negative values.

Then

\[
X=a
\]

and

\[
X=-a
\]

both produce

\[
Y=a^2.
\]

Therefore,

\[
\boxed{
P(Y=a^2)
=
P(X=a)+P(X=-a)
}
\]

in a discrete setting.

%%%%%%%%%%%%%%%%%%%%%%%%%%%%%%%%%%%%%%%%%%%%%%%%%%%%%%%%%%%%
\subsection{Worked Example: Absolute Value}
\label{subsec:worked-absolute-value-rv}
%%%%%%%%%%%%%%%%%%%%%%%%%%%%%%%%%%%%%%%%%%%%%%%%%%%%%%%%%%%%

\begin{workedexamplebox}{Many-to-One Mapping}

Suppose \(X\) satisfies

\[
P(X=-2)=0.2,
\]

\[
P(X=-1)=0.3,
\]

\[
P(X=1)=0.1,
\]

and

\[
P(X=2)=0.4.
\]

Define

\[
Y=|X|.
\]

Then

\[
P(Y=1)
=
P(X=-1)+P(X=1)
=
0.4,
\]

and

\[
P(Y=2)
=
P(X=-2)+P(X=2)
=
0.6.
\]

\begin{answerbox}
\[
\boxed{
P(Y=1)=0.4,
\qquad
P(Y=2)=0.6.
}
\]
\end{answerbox}

\end{workedexamplebox}

%%%%%%%%%%%%%%%%%%%%%%%%%%%%%%%%%%%%%%%%%%%%%%%%%%%%%%%%%%%%
\subsection{Transformation of a Discrete PMF}
\label{subsec:discrete-transformation-pmf}
%%%%%%%%%%%%%%%%%%%%%%%%%%%%%%%%%%%%%%%%%%%%%%%%%%%%%%%%%%%%

If

\[
Y=g(X),
\]

then for a discrete random variable,

\[
\boxed{
P(Y=y)
=
\sum_{\{x:g(x)=y\}}
P(X=x).
}
\]

Thus one adds the probabilities of every original value mapping to
\(y\).

%%%%%%%%%%%%%%%%%%%%%%%%%%%%%%%%%%%%%%%%%%%%%%%%%%%%%%%%%%%%
\subsection{Transforming the Support}
\label{subsec:transforming-support}
%%%%%%%%%%%%%%%%%%%%%%%%%%%%%%%%%%%%%%%%%%%%%%%%%%%%%%%%%%%%

If

\[
Y=g(X),
\]

then the possible values of \(Y\) belong to

\[
\boxed{
g(\operatorname{supp}(X)).
}
\]

Before calculating probabilities, it is often helpful to determine the
transformed support.

This prevents assigning probability to impossible values.

%%%%%%%%%%%%%%%%%%%%%%%%%%%%%%%%%%%%%%%%%%%%%%%%%%%%%%%%%%%%
\subsection{Linear Transformations}
\label{subsec:linear-transformations-rv}
%%%%%%%%%%%%%%%%%%%%%%%%%%%%%%%%%%%%%%%%%%%%%%%%%%%%%%%%%%%%

A particularly important transformation is

\[
\boxed{
Y=aX+b.
}
\]

Here:

\[
a
\]

controls scale and possibly reflection, while

\[
b
\]

controls translation.

If \(a>0\), larger values of \(X\) remain larger values of \(Y\).

If \(a<0\), the ordering reverses.

%%%%%%%%%%%%%%%%%%%%%%%%%%%%%%%%%%%%%%%%%%%%%%%%%%%%%%%%%%%%
\subsection{Location Transformations}
\label{subsec:location-transformations}
%%%%%%%%%%%%%%%%%%%%%%%%%%%%%%%%%%%%%%%%%%%%%%%%%%%%%%%%%%%%

A transformation

\[
Y=X+b
\]

shifts every possible value by \(b\).

Thus,

\[
\boxed{
\operatorname{supp}(Y)
=
\operatorname{supp}(X)+b.
}
\]

The shape of the distribution is unchanged; only its location moves.

%%%%%%%%%%%%%%%%%%%%%%%%%%%%%%%%%%%%%%%%%%%%%%%%%%%%%%%%%%%%
\subsection{Scale Transformations}
\label{subsec:scale-transformations}
%%%%%%%%%%%%%%%%%%%%%%%%%%%%%%%%%%%%%%%%%%%%%%%%%%%%%%%%%%%%

A transformation

\[
Y=aX
\]

rescales the random variable.

If

\[
a>1,
\]

distances between values increase.

If

\[
0<a<1,
\]

they decrease.

If

\[
a<0,
\]

the distribution is also reflected.

%%%%%%%%%%%%%%%%%%%%%%%%%%%%%%%%%%%%%%%%%%%%%%%%%%%%%%%%%%%%
\subsection{CDF Under Increasing Transformations}
\label{subsec:cdf-increasing-transform}
%%%%%%%%%%%%%%%%%%%%%%%%%%%%%%%%%%%%%%%%%%%%%%%%%%%%%%%%%%%%

Suppose

\[
Y=g(X)
\]

and \(g\) is strictly increasing.

Then

\[
Y\leq y
\]

is equivalent to

\[
X\leq g^{-1}(y).
\]

Therefore,

\[
\boxed{
F_Y(y)
=
F_X(g^{-1}(y)).
}
\]

This method is often called the CDF transformation method.

%%%%%%%%%%%%%%%%%%%%%%%%%%%%%%%%%%%%%%%%%%%%%%%%%%%%%%%%%%%%
\subsection{CDF Under Decreasing Transformations}
\label{subsec:cdf-decreasing-transform}
%%%%%%%%%%%%%%%%%%%%%%%%%%%%%%%%%%%%%%%%%%%%%%%%%%%%%%%%%%%%

If \(g\) is strictly decreasing, inequalities reverse.

Then

\[
g(X)\leq y
\]

corresponds to

\[
X\geq g^{-1}(y).
\]

Thus complement probabilities are often needed.

This is a common source of sign errors when transforming random
variables.

%%%%%%%%%%%%%%%%%%%%%%%%%%%%%%%%%%%%%%%%%%%%%%%%%%%%%%%%%%%%
\subsection{Worked Example: Linear CDF Transformation}
\label{subsec:worked-linear-cdf-transform}
%%%%%%%%%%%%%%%%%%%%%%%%%%%%%%%%%%%%%%%%%%%%%%%%%%%%%%%%%%%%

\begin{workedexamplebox}{Scaling and Shifting}

Let

\[
Y=2X+3.
\]

Since the transformation is increasing,

\[
Y\leq y
\]

is equivalent to

\[
2X+3\leq y.
\]

Therefore,

\[
X\leq\frac{y-3}{2}.
\]

Hence,

\begin{answerbox}
\[
\boxed{
F_Y(y)
=
F_X\left(
\frac{y-3}{2}
\right).
}
\]
\end{answerbox}

\end{workedexamplebox}

%%%%%%%%%%%%%%%%%%%%%%%%%%%%%%%%%%%%%%%%%%%%%%%%%%%%%%%%%%%%
\subsection{Continuous Transformation Preview}
\label{subsec:continuous-transformation-preview}
%%%%%%%%%%%%%%%%%%%%%%%%%%%%%%%%%%%%%%%%%%%%%%%%%%%%%%%%%%%%

If \(X\) has density \(f_X\) and

\[
Y=g(X)
\]

for a differentiable one-to-one transformation, then the density of
\(Y\) involves the derivative of the inverse transformation.

For a monotone \(g\),

\[
\boxed{
f_Y(y)
=
f_X(g^{-1}(y))
\left|
\frac{d}{dy}g^{-1}(y)
\right|.
}
\]

This change-of-variables formula will be developed more fully in
Section~\ref{sec:functions-random-variables}.

%%%%%%%%%%%%%%%%%%%%%%%%%%%%%%%%%%%%%%%%%%%%%%%%%%%%%%%%%%%%
\subsection{Why the Absolute Value Appears}
\label{subsec:jacobian-absolute-value-preview}
%%%%%%%%%%%%%%%%%%%%%%%%%%%%%%%%%%%%%%%%%%%%%%%%%%%%%%%%%%%%

The factor

\[
\left|
\frac{d}{dy}g^{-1}(y)
\right|
\]

corrects for stretching or compression of intervals.

The absolute value ensures that density remains nonnegative even when
the transformation reverses orientation.

This is the one-dimensional version of the Jacobian method used later
for multivariable transformations.

%%%%%%%%%%%%%%%%%%%%%%%%%%%%%%%%%%%%%%%%%%%%%%%%%%%%%%%%%%%%
\subsection{Quantiles}
\label{subsec:random-variable-quantiles}
%%%%%%%%%%%%%%%%%%%%%%%%%%%%%%%%%%%%%%%%%%%%%%%%%%%%%%%%%%%%

Quantiles describe positions within a probability distribution.

Informally, a \(p\)-quantile is a value \(q_p\) satisfying

\[
\boxed{
P(X\leq q_p)
\geq p
}
\]

while values below \(q_p\) have accumulated probability no greater than
\(p\).

A common formal definition is

\[
\boxed{
q_p
=
\inf
\{
x:F_X(x)\geq p
\}.
}
\]

The symbol

\[
\inf
\]

means infimum.

%%%%%%%%%%%%%%%%%%%%%%%%%%%%%%%%%%%%%%%%%%%%%%%%%%%%%%%%%%%%
\subsection{Median}
\label{subsec:random-variable-median}
%%%%%%%%%%%%%%%%%%%%%%%%%%%%%%%%%%%%%%%%%%%%%%%%%%%%%%%%%%%%

A median is a central value of the distribution.

One common quantile definition is

\[
\boxed{
m
=
q_{0.5}.
}
\]

Thus approximately half of the probability lies at or below the median
and half lies at or above it.

For discrete distributions, a median need not be unique under every
equivalent definition.

%%%%%%%%%%%%%%%%%%%%%%%%%%%%%%%%%%%%%%%%%%%%%%%%%%%%%%%%%%%%
\subsection{Quartiles}
\label{subsec:random-variable-quartiles}
%%%%%%%%%%%%%%%%%%%%%%%%%%%%%%%%%%%%%%%%%%%%%%%%%%%%%%%%%%%%

Important quantiles include:

\[
\boxed{
Q_1=q_{0.25},
}
\]

\[
\boxed{
Q_2=q_{0.50},
}
\]

and

\[
\boxed{
Q_3=q_{0.75}.
}
\]

These are the first, second, and third quartiles.

The second quartile is the median.

%%%%%%%%%%%%%%%%%%%%%%%%%%%%%%%%%%%%%%%%%%%%%%%%%%%%%%%%%%%%
\subsection{Percentiles}
\label{subsec:random-variable-percentiles}
%%%%%%%%%%%%%%%%%%%%%%%%%%%%%%%%%%%%%%%%%%%%%%%%%%%%%%%%%%%%

The \(100p\)-th percentile corresponds to the \(p\)-quantile.

For example:

\[
q_{0.90}
\]

is the \(90\)-th percentile.

Quantiles provide distribution summaries without requiring a specific
parametric model.

%%%%%%%%%%%%%%%%%%%%%%%%%%%%%%%%%%%%%%%%%%%%%%%%%%%%%%%%%%%%
\subsection{Inverse CDF}
\label{subsec:inverse-cdf}
%%%%%%%%%%%%%%%%%%%%%%%%%%%%%%%%%%%%%%%%%%%%%%%%%%%%%%%%%%%%

When \(F_X\) is strictly increasing and continuous, one may solve

\[
F_X(x)=p
\]

for \(x\).

Then

\[
\boxed{
x=F_X^{-1}(p).
}
\]

More generally, the quantile function is defined by a generalized
inverse:

\[
\boxed{
F_X^{-1}(p)
=
\inf
\{
x:F_X(x)\geq p
\}.
}
\]

%%%%%%%%%%%%%%%%%%%%%%%%%%%%%%%%%%%%%%%%%%%%%%%%%%%%%%%%%%%%
\subsection{Random Variable Notation}
\label{subsec:random-variable-notation}
%%%%%%%%%%%%%%%%%%%%%%%%%%%%%%%%%%%%%%%%%%%%%%%%%%%%%%%%%%%%

Random variables are conventionally written using capital letters:

\[
X,
\quad
Y,
\quad
Z.
\]

Realized numerical values are often written using lowercase letters:

\[
x,
\quad
y,
\quad
z.
\]

Thus,

\[
\boxed{
P(X=x)
}
\]

means the probability that random variable \(X\) takes the particular
value \(x\).

%%%%%%%%%%%%%%%%%%%%%%%%%%%%%%%%%%%%%%%%%%%%%%%%%%%%%%%%%%%%
\subsection{Distribution Notation}
\label{subsec:distribution-notation}
%%%%%%%%%%%%%%%%%%%%%%%%%%%%%%%%%%%%%%%%%%%%%%%%%%%%%%%%%%%%

One may write

\[
\boxed{
X\sim\mathcal{D}
}
\]

to indicate that \(X\) follows distribution \(\mathcal{D}\).

The symbol

\[
\sim
\]

is read here as

\[
\text{``is distributed as.''}
\]

For example, later we will write expressions such as

\[
X\sim\operatorname{Bin}(n,p).
\]

%%%%%%%%%%%%%%%%%%%%%%%%%%%%%%%%%%%%%%%%%%%%%%%%%%%%%%%%%%%%
\subsection{Equality in Distribution}
\label{subsec:equality-in-distribution}
%%%%%%%%%%%%%%%%%%%%%%%%%%%%%%%%%%%%%%%%%%%%%%%%%%%%%%%%%%%%

Two random variables \(X\) and \(Y\) are equal in distribution if they
have the same probability law.

One notation is

\[
\boxed{
X\overset{d}{=}Y.
}
\]

This means

\[
\boxed{
F_X(x)=F_Y(x)
}
\]

for every real \(x\).

The variables need not be defined on the same probability space.

%%%%%%%%%%%%%%%%%%%%%%%%%%%%%%%%%%%%%%%%%%%%%%%%%%%%%%%%%%%%
\subsection{Equality Almost Surely Preview}
\label{subsec:equality-almost-surely-preview}
%%%%%%%%%%%%%%%%%%%%%%%%%%%%%%%%%%%%%%%%%%%%%%%%%%%%%%%%%%%%

Another important concept is equality with probability \(1\).

We write

\[
\boxed{
X=Y
\quad\text{almost surely}
}
\]

if

\[
\boxed{
P(X=Y)=1.
}
\]

This is stronger than equality in distribution.

The abbreviation

\[
\text{a.s.}
\]

stands for ``almost surely.''

The measure-theoretic meaning will be developed later.

%%%%%%%%%%%%%%%%%%%%%%%%%%%%%%%%%%%%%%%%%%%%%%%%%%%%%%%%%%%%
\subsection{Equality in Distribution Is Not Pointwise Equality}
\label{subsec:distribution-not-pointwise}
%%%%%%%%%%%%%%%%%%%%%%%%%%%%%%%%%%%%%%%%%%%%%%%%%%%%%%%%%%%%

If

\[
X\overset{d}{=}Y,
\]

then \(X\) and \(Y\) have the same distribution.

This does not imply

\[
X(\omega)=Y(\omega)
\]

for each outcome.

Thus:

\[
\boxed{
\text{same distribution}
\neq
\text{same random variable pointwise}.
}
\]

%%%%%%%%%%%%%%%%%%%%%%%%%%%%%%%%%%%%%%%%%%%%%%%%%%%%%%%%%%%%
\subsection{Random Vectors Preview}
\label{subsec:random-vectors-preview}
%%%%%%%%%%%%%%%%%%%%%%%%%%%%%%%%%%%%%%%%%%%%%%%%%%%%%%%%%%%%

Several random variables may be grouped together.

A random vector has the form

\[
\boxed{
\mathbf{X}
=
(X_1,\ldots,X_n).
}
\]

It maps outcomes into

\[
\mathbb{R}^n.
\]

Thus,

\[
\boxed{
\mathbf{X}:\Omega\to\mathbb{R}^n.
}
\]

Random vectors are the foundation of joint distributions.

%%%%%%%%%%%%%%%%%%%%%%%%%%%%%%%%%%%%%%%%%%%%%%%%%%%%%%%%%%%%
\subsection{Bivariate Random Variables}
\label{subsec:bivariate-random-variables}
%%%%%%%%%%%%%%%%%%%%%%%%%%%%%%%%%%%%%%%%%%%%%%%%%%%%%%%%%%%%

A pair

\[
(X,Y)
\]

is called a bivariate random variable or random vector.

For example, when rolling two dice, define

\[
X
=
\text{first die},
\]

and

\[
Y
=
\text{second die}.
\]

Then

\[
(X,Y)
\]

records both numerical outcomes simultaneously.

Their dependence structure will be developed in
Section~\ref{sec:joint-distributions}.

%%%%%%%%%%%%%%%%%%%%%%%%%%%%%%%%%%%%%%%%%%%%%%%%%%%%%%%%%%%%
\subsection{Random Variables from Measurements}
\label{subsec:rv-measurements}
%%%%%%%%%%%%%%%%%%%%%%%%%%%%%%%%%%%%%%%%%%%%%%%%%%%%%%%%%%%%

In applications, random variables often represent measurements.

Examples include:

\[
X
=
\text{temperature},
\]

\[
Y
=
\text{aircraft delay time},
\]

\[
Z
=
\text{dissolved oxygen concentration},
\]

or

\[
N
=
\text{number of arrivals}.
\]

The mathematical distribution represents uncertainty or variability in
the observed quantity.

%%%%%%%%%%%%%%%%%%%%%%%%%%%%%%%%%%%%%%%%%%%%%%%%%%%%%%%%%%%%
\subsection{Random Variables from Counts}
\label{subsec:rv-counts}
%%%%%%%%%%%%%%%%%%%%%%%%%%%%%%%%%%%%%%%%%%%%%%%%%%%%%%%%%%%%

Many discrete random variables count occurrences.

Examples include:

\[
X
=
\text{number of defective items},
\]

\[
Y
=
\text{number of successes},
\]

\[
Z
=
\text{number of graph edges},
\]

and

\[
N
=
\text{number of customers}.
\]

Count variables naturally take values in

\[
\boxed{
\mathbb{Z}_{\geq0}
=
\{0,1,2,\ldots\}.
}
\]

%%%%%%%%%%%%%%%%%%%%%%%%%%%%%%%%%%%%%%%%%%%%%%%%%%%%%%%%%%%%
\subsection{Random Variables from Waiting Times}
\label{subsec:rv-waiting-times}
%%%%%%%%%%%%%%%%%%%%%%%%%%%%%%%%%%%%%%%%%%%%%%%%%%%%%%%%%%%%

Waiting times may be either discrete or continuous.

For example,

\[
X
=
\text{number of trials until the first success}
\]

is discrete.

By contrast,

\[
T
=
\text{physical time until a component fails}
\]

may be continuous.

The mathematical model depends on how the experiment is measured.

%%%%%%%%%%%%%%%%%%%%%%%%%%%%%%%%%%%%%%%%%%%%%%%%%%%%%%%%%%%%
\subsection{Random Variables from Indicators}
\label{subsec:rv-from-indicators}
%%%%%%%%%%%%%%%%%%%%%%%%%%%%%%%%%%%%%%%%%%%%%%%%%%%%%%%%%%%%

A complicated count can often be decomposed into simple indicators.

For example, suppose a graph has potential edges

\[
e_1,\ldots,e_m.
\]

Let

\[
I_j
=
\mathbf{1}_{\{e_j\text{ appears}\}}.
\]

Then the total number of edges is

\[
\boxed{
X
=
\sum_{j=1}^{m}I_j.
}
\]

This representation becomes especially powerful when computing
expectations.

%%%%%%%%%%%%%%%%%%%%%%%%%%%%%%%%%%%%%%%%%%%%%%%%%%%%%%%%%%%%
\subsection{Random Variables and Information Reduction}
\label{subsec:rv-information-reduction}
%%%%%%%%%%%%%%%%%%%%%%%%%%%%%%%%%%%%%%%%%%%%%%%%%%%%%%%%%%%%

A random variable often records only one feature of an experiment.

For example, a five-card poker hand contains substantial information.

The random variable

\[
X
=
\text{number of aces}
\]

retains only one statistic.

Thus a random variable may be viewed as a numerical summary:

\[
\boxed{
\text{complex random outcome}
\longrightarrow
\text{numerical feature}.
}
\]

%%%%%%%%%%%%%%%%%%%%%%%%%%%%%%%%%%%%%%%%%%%%%%%%%%%%%%%%%%%%
\subsection{Statistics as Random Variables}
\label{subsec:statistics-as-random-variables}
%%%%%%%%%%%%%%%%%%%%%%%%%%%%%%%%%%%%%%%%%%%%%%%%%%%%%%%%%%%%

Before data are observed, a statistic computed from a random sample is
itself a random variable.

For example, if

\[
X_1,\ldots,X_n
\]

are random observations, then the sample mean

\[
\overline{X}
=
\frac1n
\sum_{i=1}^{n}X_i
\]

is a random variable.

This idea becomes fundamental in statistical inference.

%%%%%%%%%%%%%%%%%%%%%%%%%%%%%%%%%%%%%%%%%%%%%%%%%%%%%%%%%%%%
\subsection{Random Variables and Measurability Preview}
\label{subsec:rv-measurability-preview}
%%%%%%%%%%%%%%%%%%%%%%%%%%%%%%%%%%%%%%%%%%%%%%%%%%%%%%%%%%%%

In elementary finite probability, nearly any numerical function on the
sample space can be used as a random variable.

In general probability spaces, a random variable must satisfy a
measurability condition.

Informally, events such as

\[
\{X\leq x\}
\]

must belong to the collection of allowable events.

The formal definition will appear in measure-theoretic probability.

%%%%%%%%%%%%%%%%%%%%%%%%%%%%%%%%%%%%%%%%%%%%%%%%%%%%%%%%%%%%
\subsection{Distribution Functions Determine Laws}
\label{subsec:cdf-determines-law}
%%%%%%%%%%%%%%%%%%%%%%%%%%%%%%%%%%%%%%%%%%%%%%%%%%%%%%%%%%%%

A CDF contains complete information about a real-valued probability
distribution.

If

\[
F_X(x)=F_Y(x)
\]

for every \(x\), then

\[
\boxed{
X\overset{d}{=}Y.
}
\]

Therefore the distribution of a real random variable may be specified
entirely by its CDF.

%%%%%%%%%%%%%%%%%%%%%%%%%%%%%%%%%%%%%%%%%%%%%%%%%%%%%%%%%%%%
\subsection{Validating a Proposed CDF}
\label{subsec:validating-cdf}
%%%%%%%%%%%%%%%%%%%%%%%%%%%%%%%%%%%%%%%%%%%%%%%%%%%%%%%%%%%%

To determine whether a function \(F\) can be a CDF, verify:

\begin{enumerate}
    \item
    \[
    F
    \text{ is nondecreasing};
    \]

    \item
    \[
    \lim_{x\to-\infty}F(x)=0;
    \]

    \item
    \[
    \lim_{x\to\infty}F(x)=1;
    \]

    \item \(F\) is right-continuous.
\end{enumerate}

These conditions characterize CDFs on the real line.

%%%%%%%%%%%%%%%%%%%%%%%%%%%%%%%%%%%%%%%%%%%%%%%%%%%%%%%%%%%%
\subsection{Worked Example: Is This a CDF?}
\label{subsec:worked-valid-cdf}
%%%%%%%%%%%%%%%%%%%%%%%%%%%%%%%%%%%%%%%%%%%%%%%%%%%%%%%%%%%%

\begin{workedexamplebox}{Checking a Candidate}

Consider

\[
F(x)
=
\begin{cases}
0,&x<0,\\
x,&0\leq x\leq1,\\
1,&x>1.
\end{cases}
\]

The function is nondecreasing.

Also,

\[
\lim_{x\to-\infty}F(x)=0
\]

and

\[
\lim_{x\to\infty}F(x)=1.
\]

It is continuous, and therefore right-continuous.

\begin{answerbox}
\[
\boxed{
F
\text{ is a valid CDF}.
}
\]
\end{answerbox}

\end{workedexamplebox}

%%%%%%%%%%%%%%%%%%%%%%%%%%%%%%%%%%%%%%%%%%%%%%%%%%%%%%%%%%%%
\subsection{Distribution Functions with Jumps}
\label{subsec:cdf-with-jumps}
%%%%%%%%%%%%%%%%%%%%%%%%%%%%%%%%%%%%%%%%%%%%%%%%%%%%%%%%%%%%

A CDF may contain jumps.

For example,

\[
F(x)
=
\begin{cases}
0,&x<0,\\
0.3,&0\leq x<1,\\
1,&x\geq1.
\end{cases}
\]

has jump sizes

\[
0.3
\]

at \(0\) and

\[
0.7
\]

at \(1\).

Therefore,

\[
P(X=0)=0.3
\]

and

\[
P(X=1)=0.7.
\]

%%%%%%%%%%%%%%%%%%%%%%%%%%%%%%%%%%%%%%%%%%%%%%%%%%%%%%%%%%%%
\subsection{CDFs Need Not Be Differentiable}
\label{subsec:cdf-not-differentiable}
%%%%%%%%%%%%%%%%%%%%%%%%%%%%%%%%%%%%%%%%%%%%%%%%%%%%%%%%%%%%

Every random variable has a CDF.

Not every CDF has a derivative everywhere.

For discrete variables, the CDF has jumps.

For mixed variables, it may contain both smooth regions and jumps.

Thus the CDF is more general than either a PMF or PDF.

%%%%%%%%%%%%%%%%%%%%%%%%%%%%%%%%%%%%%%%%%%%%%%%%%%%%%%%%%%%%
\subsection{PMF, PDF, and CDF Comparison}
\label{subsec:pmf-pdf-cdf-comparison}
%%%%%%%%%%%%%%%%%%%%%%%%%%%%%%%%%%%%%%%%%%%%%%%%%%%%%%%%%%%%

\[
\begin{array}{c|c|c}
\text{Object}
&
\text{Typical Use}
&
\text{Probability Calculation}
\\
\hline
p_X(x)
&
\text{Discrete}
&
P(X=x)=p_X(x)
\\[1mm]
f_X(x)
&
\text{Continuous}
&
P(a<X<b)=\displaystyle\int_a^b f_X(x)\,dx
\\[3mm]
F_X(x)
&
\text{General}
&
P(X\leq x)=F_X(x)
\end{array}
\]

The CDF applies to all real-valued random variables.

%%%%%%%%%%%%%%%%%%%%%%%%%%%%%%%%%%%%%%%%%%%%%%%%%%%%%%%%%%%%
\subsection{Common Random-Variable Workflow}
\label{subsec:random-variable-workflow}
%%%%%%%%%%%%%%%%%%%%%%%%%%%%%%%%%%%%%%%%%%%%%%%%%%%%%%%%%%%%

When analyzing a random variable:

\begin{enumerate}
    \item identify the underlying experiment;
    \item define the sample space if necessary;
    \item specify the function \(X(\omega)\);
    \item determine the possible values of \(X\);
    \item classify \(X\) as discrete, continuous, or mixed;
    \item identify its PMF, PDF, or CDF;
    \item verify normalization;
    \item translate requested conditions into events involving \(X\);
    \item compute the desired probability;
    \item check that the result lies in \([0,1]\).
\end{enumerate}

%%%%%%%%%%%%%%%%%%%%%%%%%%%%%%%%%%%%%%%%%%%%%%%%%%%%%%%%%%%%
\subsection{Common Errors}
\label{subsec:random-variable-common-errors}
%%%%%%%%%%%%%%%%%%%%%%%%%%%%%%%%%%%%%%%%%%%%%%%%%%%%%%%%%%%%

\subsubsection{Treating a Random Variable as an Outcome}

A random variable is a function on outcomes.

It is not itself one outcome.

%%%%%%%%%%%%%%%%%%%%%%%%%%%%%%%%%%%%%%%%%%%%%%%%%%%%%%%%%%%%
\subsubsection{Confusing \(X\) with \(x\)}

Conventionally,

\[
X
\]

denotes the random variable, while

\[
x
\]

denotes one possible numerical value.

%%%%%%%%%%%%%%%%%%%%%%%%%%%%%%%%%%%%%%%%%%%%%%%%%%%%%%%%%%%%
\subsubsection{Confusing PMF and PDF}

For a discrete random variable,

\[
P(X=x)
=
p_X(x).
\]

For a continuous random variable,

\[
P(X=x)=0
\]

even when

\[
f_X(x)>0.
\]

%%%%%%%%%%%%%%%%%%%%%%%%%%%%%%%%%%%%%%%%%%%%%%%%%%%%%%%%%%%%
\subsubsection{Using a Density Value as a Probability}

The quantity

\[
f_X(x)
\]

may exceed \(1\).

This does not violate probability laws.

Only integrated probability must lie between \(0\) and \(1\).

%%%%%%%%%%%%%%%%%%%%%%%%%%%%%%%%%%%%%%%%%%%%%%%%%%%%%%%%%%%%
\subsubsection{Forgetting PMF Normalization}

A valid discrete distribution requires

\[
\boxed{
\sum_xp_X(x)=1.
}
\]

%%%%%%%%%%%%%%%%%%%%%%%%%%%%%%%%%%%%%%%%%%%%%%%%%%%%%%%%%%%%
\subsubsection{Forgetting PDF Normalization}

A valid density requires

\[
\boxed{
\int_{-\infty}^{\infty}f_X(x)\,dx=1.
}
\]

%%%%%%%%%%%%%%%%%%%%%%%%%%%%%%%%%%%%%%%%%%%%%%%%%%%%%%%%%%%%
\subsubsection{Assuming Every Random Variable Is Discrete or Continuous}

Mixed distributions also exist.

%%%%%%%%%%%%%%%%%%%%%%%%%%%%%%%%%%%%%%%%%%%%%%%%%%%%%%%%%%%%
\subsubsection{Forgetting Endpoint Masses}

For discrete variables,

\[
P(a<X\leq b)
\]

and

\[
P(a\leq X\leq b)
\]

may differ.

Endpoint inclusion matters whenever there is positive probability mass
at an endpoint.

%%%%%%%%%%%%%%%%%%%%%%%%%%%%%%%%%%%%%%%%%%%%%%%%%%%%%%%%%%%%
\subsubsection{Assuming Every CDF Is Continuous}

Discrete random variables have CDFs with jumps.

A CDF is required to be right-continuous, not necessarily continuous.

%%%%%%%%%%%%%%%%%%%%%%%%%%%%%%%%%%%%%%%%%%%%%%%%%%%%%%%%%%%%
\subsubsection{Ignoring Multiple Preimages Under Transformations}

If

\[
Y=g(X)
\]

and several \(x\)-values produce the same \(y\), all corresponding
probabilities must be combined.

%%%%%%%%%%%%%%%%%%%%%%%%%%%%%%%%%%%%%%%%%%%%%%%%%%%%%%%%%%%%
\subsubsection{Forgetting to Reverse Inequalities}

If

\[
Y=aX+b
\]

with

\[
a<0,
\]

solving inequalities involving \(Y\) reverses the inequality direction.

%%%%%%%%%%%%%%%%%%%%%%%%%%%%%%%%%%%%%%%%%%%%%%%%%%%%%%%%%%%%
\subsubsection{Confusing Equality in Distribution with Equality}

Two variables can have the same distribution without being the same
function or taking equal values on each outcome.

%%%%%%%%%%%%%%%%%%%%%%%%%%%%%%%%%%%%%%%%%%%%%%%%%%%%%%%%%%%%
\subsection{Applications}
\label{subsec:random-variable-applications}
%%%%%%%%%%%%%%%%%%%%%%%%%%%%%%%%%%%%%%%%%%%%%%%%%%%%%%%%%%%%

\begin{applicationbox}

\textbf{Statistics.}
Sample means, sample proportions, estimators, and test statistics are
random variables before data are observed.

\medskip

\textbf{Engineering.}
Random variables model lifetimes, loads, measurement errors, component
counts, and system performance.

\medskip

\textbf{Finance.}
Returns, losses, prices, and default indicators are represented as
random variables.

\medskip

\textbf{Biology.}
Population counts, genetic outcomes, concentrations, and survival times
are modeled probabilistically.

\medskip

\textbf{Computer Science.}
Algorithm runtimes, hash collisions, random graph statistics, and
randomized algorithm outputs are random variables.

\medskip

\textbf{Operations Research.}
Demand, service time, waiting time, inventory usage, and arrival counts
are modeled as random variables.

\medskip

\textbf{Physics.}
Measured positions, energies, particle counts, and fluctuations can be
represented probabilistically.

\medskip

\textbf{Combinatorics.}
Counts of edges, fixed points, cycles, inversions, and substructures in
random discrete objects are random variables.

\end{applicationbox}

%%%%%%%%%%%%%%%%%%%%%%%%%%%%%%%%%%%%%%%%%%%%%%%%%%%%%%%%%%%%
\subsection{Exercises}
\label{subsec:random-variable-exercises}
%%%%%%%%%%%%%%%%%%%%%%%%%%%%%%%%%%%%%%%%%%%%%%%%%%%%%%%%%%%%

\begin{exercise}[1]
Define a random variable.
\end{exercise}

\begin{exercise}[1]
Explain why a random variable is a function.
\end{exercise}

\begin{exercise}[1]
For two coin tosses, define \(X\) to be the number of tails and list all
values of \(X\).
\end{exercise}

\begin{exercise}[1]
For one die roll, define

\[
X=2(\text{die result})+1.
\]

Find the support of \(X\).
\end{exercise}

\begin{exercise}[1]
Define a discrete random variable.
\end{exercise}

\begin{exercise}[1]
Define a probability mass function.
\end{exercise}

\begin{exercise}[1]
State the two basic PMF conditions.
\end{exercise}

\begin{exercise}[1]
Define the cumulative distribution function.
\end{exercise}

\begin{exercise}[1]
State four fundamental properties of a CDF.
\end{exercise}

\begin{exercise}[1]
Define an indicator random variable.
\end{exercise}

\begin{exercise}[1]
What is the support of an indicator variable?
\end{exercise}

\begin{exercise}[1]
Define a degenerate random variable.
\end{exercise}

\begin{exercise}[2]
A random variable \(X\) has PMF

\[
p_X(0)=0.2,
\qquad
p_X(1)=0.5,
\qquad
p_X(2)=0.3.
\]

Find

\[
P(X\leq1).
\]
\end{exercise}

\begin{exercise}[2]
Using the previous PMF, find

\[
P(X>0).
\]
\end{exercise}

\begin{exercise}[2]
Construct the complete CDF of the previous random variable.
\end{exercise}

\begin{exercise}[2]
Suppose

\[
P(X=-1)=0.25,
\]

\[
P(X=0)=0.50,
\]

and

\[
P(X=1)=0.25.
\]

Define

\[
Y=X^2.
\]

Find the PMF of \(Y\).
\end{exercise}

\begin{exercise}[2]
Let

\[
Y=3X-2.
\]

If

\[
\operatorname{supp}(X)=\{0,1,2,3\},
\]

find the support of \(Y\).
\end{exercise}

\begin{exercise}[2]
Suppose

\[
F_X(2)=0.4
\]

and

\[
F_X(7)=0.9.
\]

Find

\[
P(2<X\leq7).
\]
\end{exercise}

\begin{exercise}[2]
Suppose a CDF jumps from \(0.35\) to \(0.60\) at \(x=4\).

Find

\[
P(X=4).
\]
\end{exercise}

\begin{exercise}[2]
A continuous random variable has density

\[
f(x)
=
\begin{cases}
\dfrac12,&0\leq x\leq2,\\
0,&\text{otherwise}.
\end{cases}
\]

Find

\[
P(0.5<X<1.5).
\]
\end{exercise}

\begin{exercise}[2]
Find the CDF of the random variable in the previous exercise.
\end{exercise}

\begin{exercise}[2]
Explain why

\[
P(X=1)=0
\]

for the continuous random variable above.
\end{exercise}

\begin{exercise}[2]
Suppose

\[
F_X(x)
=
\begin{cases}
0,&x<0,\\
x^2,&0\leq x\leq1,\\
1,&x>1.
\end{cases}
\]

Find

\[
P\left(
\frac14<X\leq\frac34
\right).
\]
\end{exercise}

\begin{exercise}[2]
Find the density corresponding to the CDF in the previous exercise.
\end{exercise}

\begin{exercise}[2]
Let

\[
I=\mathbf{1}_A
\]

where

\[
P(A)=0.3.
\]

Find the PMF of \(I\).
\end{exercise}

\begin{exercise}[3]
Prove that every CDF is nondecreasing.
\end{exercise}

\begin{exercise}[3]
Show that

\[
P(a<X\leq b)
=
F_X(b)-F_X(a).
\]
\end{exercise}

\begin{exercise}[3]
Show that

\[
P(X=x)
=
F_X(x)-F_X(x^-).
\]
\end{exercise}

\begin{exercise}[3]
Prove that a continuous random variable satisfies

\[
P(X=x)=0
\]

for every \(x\).
\end{exercise}

\begin{exercise}[3]
A discrete random variable has PMF

\[
p_X(k)=c2^{-k},
\qquad
k=1,2,3,\ldots.
\]

Determine \(c\).
\end{exercise}

\begin{exercise}[3]
Find the CDF of the random variable in the previous exercise.
\end{exercise}

\begin{exercise}[3]
Let \(X\) have support

\[
\{-2,-1,0,1,2\}
\]

with equal probabilities.

Find the distribution of

\[
Y=|X|.
\]
\end{exercise}

\begin{exercise}[3]
Using the previous random variable, find the distribution of

\[
Z=X^2.
\]
\end{exercise}

\begin{exercise}[3]
Suppose

\[
Y=5-2X.
\]

Express

\[
F_Y(y)
\]

in terms of \(F_X\), carefully accounting for the reversed inequality.
\end{exercise}

\begin{exercise}[3]
Suppose \(X\) is continuous with density \(f_X\), and

\[
Y=2X+1.
\]

Use the change-of-variables formula to express \(f_Y\) in terms of
\(f_X\).
\end{exercise}

\begin{exercise}[3]
Prove that if

\[
X\overset{d}{=}Y,
\]

then their CDFs are equal.
\end{exercise}

\begin{exercise}[3]
Give an example of two random variables that have the same distribution
but are not equal pointwise.
\end{exercise}

\begin{exercise}[3]
For a random graph on \(n\) vertices, define indicator variables for
each possible edge and express the total number of edges as their sum.
\end{exercise}

\begin{exercise}[3]
Let \(X\) be the number of fixed points of a random permutation.
Construct indicator variables whose sum equals \(X\).
\end{exercise}

\begin{exercise}[3]
Let

\[
q_p
=
\inf\{x:F_X(x)\geq p\}.
\]

Explain why this definition is useful for discrete distributions with
jumps.
\end{exercise}

\begin{exercise}[4]
Study the formal measurability requirement in the definition of a random
variable.
\end{exercise}

\begin{exercise}[4]
Investigate distributions that are neither discrete nor absolutely
continuous.
\end{exercise}

\begin{exercise}[4]
Study mixed distributions containing both point masses and densities.
\end{exercise}

\begin{exercise}[4]
Investigate the Cantor distribution as an example of a singular
continuous distribution.
\end{exercise}

\begin{exercise}[4]
Study generalized inverse distribution functions and quantile
transformations.
\end{exercise}

\begin{exercise}[4]
Prove that if

\[
U\sim\operatorname{Uniform}(0,1),
\]

then under suitable conditions,

\[
X=F^{-1}(U)
\]

has CDF \(F\).
\end{exercise}

\begin{exercise}[4]
Investigate probability integral transforms.
\end{exercise}

\begin{exercise}[4]
Study random vectors and the distinction between marginal and joint
distributions.
\end{exercise}

\begin{exercise}[4]
Investigate pushforward measures as the measure-theoretic definition of
the distribution of a random variable.
\end{exercise}

\begin{exercise}[4]
Study almost-sure equality and compare it with equality in distribution
and pointwise equality.
\end{exercise}

%%%%%%%%%%%%%%%%%%%%%%%%%%%%%%%%%%%%%%%%%%%%%%%%%%%%%%%%%%%%
\subsection{Conceptual Summary}
\label{subsec:random-variable-summary}
%%%%%%%%%%%%%%%%%%%%%%%%%%%%%%%%%%%%%%%%%%%%%%%%%%%%%%%%%%%%

A random variable is a function

\[
\boxed{
X:\Omega\to\mathbb{R}.
}
\]

It converts random outcomes into numerical values.

Statements such as

\[
X\leq x
\]

represent events:

\[
\boxed{
\{X\leq x\}
=
\{\omega:X(\omega)\leq x\}.
}
\]

For a discrete random variable, the probability mass function is

\[
\boxed{
p_X(x)=P(X=x).
}
\]

It satisfies

\[
\boxed{
p_X(x)\geq0,
\qquad
\sum_xp_X(x)=1.
}
\]

For a continuous random variable, probabilities are described by a
density:

\[
\boxed{
P(a<X\leq b)
=
\int_a^b f_X(x)\,dx.
}
\]

The cumulative distribution function is

\[
\boxed{
F_X(x)=P(X\leq x).
}
\]

Every real-valued random variable has a CDF.

The CDF is:

\[
\boxed{
\text{nondecreasing},
}
\]

\[
\boxed{
\text{right-continuous},
}
\]

with

\[
\boxed{
F_X(-\infty)=0,
\qquad
F_X(\infty)=1.
}
\]

For a discrete random variable,

\[
\boxed{
P(X=x)
=
F_X(x)-F_X(x^-).
}
\]

For an absolutely continuous random variable,

\[
\boxed{
F_X(x)
=
\int_{-\infty}^{x}f_X(t)\,dt.
}
\]

An indicator variable is

\[
\boxed{
\mathbf{1}_A
=
\begin{cases}
1,&A\text{ occurs},\\
0,&A\text{ does not occur}.
\end{cases}
}
\]

Functions of random variables produce new random variables:

\[
\boxed{
Y=g(X).
}
\]

The distribution of the transformed variable depends on both the
distribution of \(X\) and the mapping \(g\).

Quantiles are obtained from the CDF through

\[
\boxed{
q_p
=
\inf\{x:F_X(x)\geq p\}.
}
\]

Random vectors extend the idea to several numerical quantities:

\[
\boxed{
\mathbf{X}
=
(X_1,\ldots,X_n).
}
\]

Random variables therefore provide the fundamental bridge

\[
\boxed{
\text{probability space}
\longrightarrow
\text{numerical distributions}
\longrightarrow
\text{analysis}.
}
\]

%%%%%%%%%%%%%%%%%%%%%%%%%%%%%%%%%%%%%%%%%%%%%%%%%%%%%%%%%%%%
\subsection{Section Connections}
\label{subsec:random-variable-connections}
%%%%%%%%%%%%%%%%%%%%%%%%%%%%%%%%%%%%%%%%%%%%%%%%%%%%%%%%%%%%

\chapterconnection{
Random Variables convert the event-based probability framework of
Section~\ref{sec:elementary-probability} into numerical probability
models. The next two sections develop Discrete Distributions and
Continuous Distributions, which provide major families of random
variables and their probability laws. Joint Distributions extends the
framework to random vectors and dependence between several variables.
Expectation and Moments assign numerical summaries to distributions,
while Functions of Random Variables develops transformation techniques
systematically. Probability Transforms encode distributions through
generating and transform methods. The Laws of Large Numbers and Central
Limit Theorems study sequences of random variables, while
Measure-Theoretic Probability later provides the formal foundations for
measurability, distributions, and almost-sure statements.
}

%%%%%%%%%%%%%%%%%%%%%%%%%%%%%%%%%%%%%%%%%%%%%%%%%%%%%%%%%%%%
% SECTION SOURCE TRACKING
%%%%%%%%%%%%%%%%%%%%%%%%%%%%%%%%%%%%%%%%%%%%%%%%%%%%%%%%%%%%

%------------------------------------------------------------
% 7.2 Random Variables
%------------------------------------------------------------
%
% Source basis:
%   - Standard elementary probability theory.
%   - Standard random-variable theory.
%   - Standard introductory distribution theory.
%
% Used for:
%   - random variables as functions;
%   - outcomes versus values;
%   - preimages and random-variable events;
%   - distributions and supports;
%   - discrete random variables;
%   - probability mass functions;
%   - countably infinite support;
%   - continuous random variables;
%   - probability density functions preview;
%   - mixed distributions;
%   - degenerate random variables;
%   - cumulative distribution functions;
%   - CDF properties;
%   - discrete CDFs;
%   - CDF jumps;
%   - continuous CDFs;
%   - recovering probabilities from CDFs;
%   - PMFs and PDFs from CDFs;
%   - indicator random variables;
%   - functions of random variables;
%   - discrete transformations;
%   - transformed support;
%   - location and scale transformations;
%   - CDF transformation method;
%   - density transformation preview;
%   - quantiles;
%   - medians and quartiles;
%   - inverse CDFs;
%   - distribution notation;
%   - equality in distribution;
%   - almost-sure equality preview;
%   - random vectors preview;
%   - measurement and count variables;
%   - statistics as random variables;
%   - measurability preview;
%   - CDF validation;
%   - applications and exercises.
%
% Notes:
%   - Standard discrete families are developed in Section 7.3.
%   - Standard continuous families are developed in Section 7.4.
%   - Joint random vectors are developed in Section 7.5.
%   - Expectation is intentionally deferred to Section 7.6 except for
%     brief forward connections.
%   - Transformation theory is developed further in Section 7.7.
%   - Formal measurability and probability-law construction are
%     developed in Section 7.11.
%
%%%%%%%%%%%%%%%%%%%%%%%%%%%%%%%%%%%%%%%%%%%%%%%%%%%%%%%%%%%%